\section{Small Holes: Overview}
\label{sec:overview}

By \cref{thm:reduction}, deciding a disjunctive entailment amounts to
deciding whether a structural reader has a hole. The rest of the proof is
about readers and words; it uses no property of trees or constraints, and
the alphabet $\Alp$ is arbitrary. We fix a structural reader
$M=(S,\delta,(s_i)_i,\dead,\Rej,\Adm)$ with $k$ starts and $N$ states.

\begin{theorem}[Small holes]\label{thm:small-hole}
There is a computable function $\Hole$ such that every structural reader with
$k$ starts and $N$ states that has a hole has a hole of total length at most
$\Hole(k,N)$.
\end{theorem}

The definition of $\Hole$ is given in \cref{eq:hole-bound} below, and its
growth is estimated in \cref{sec:complexity}. This section introduces the
notions that organize the proof and reduces the theorem to two bounds.

\subsection{Domination, Fresh Words and Immortal States}
\label{sec:overview-notions}

A \emph{table} with $k$ starts is the part $(S,\delta,(s_i)_i)$ of a reader:
a complete deterministic automaton with $k$ start states. Runs, endpoints and
channels of a tuple are defined for tables as in \cref{sec:readers-def}. We
write $\Vis_i(W)=\{q_i(a)\mid 0\le a\le\len{w_i}\}$ for the states that
coordinate $i$ visits.

\begin{definition}[Domination]\label{def:domination}
A tuple $U$ \emph{dominates} a tuple $W$ in a table, written $U\dominates W$,
if for all coordinates $i$ and $j$: the endpoints of coordinate $i$ in $U$
and in $W$ are equal, $\Vis_i(U)\subseteq\Vis_i(W)$, and
$\Chan_{ij}(U)\subseteq\Chan_{ij}(W)$.
\end{definition}

Domination is reflexive and transitive. A dominating tuple does nothing that
the dominated tuple does not do: it ends in the same states, passes only
through states that were already visited, and compares only pairs of states
that were already compared. Since the definition of a hole only forbids
things, we obtain:

\begin{lemma}[Holes under domination]\label{lem:hole-dominates}
If $W$ is a hole of a structural reader and $U\dominates W$ in its table,
then $U$ is a hole with the same endpoints.
\end{lemma}

\begin{proof}
The run of $u_i$ visits only states of $\Vis_i(W)$, none of which is dead,
and its endpoint is that of $w_i$, which is rejected. No pair of
$\Chan_{ij}(U)\subseteq\Chan_{ij}(W)$ is admitted.
\end{proof}

So a hole may be replaced by any tuple that dominates it, and the task is to
find short dominating tuples. The block deletions of \cref{sec:fresh} and
the compression of periodic stretches in \cref{sec:tails-compression} delete factors from the
words of a tuple in such a way that the result dominates it. The
replacement of rejected tails by rays in \cref{sec:tails-choice} is of a
different kind: it can lengthen words and change endpoints, and it only
preserves the existence of a hole. Deleting a factor between two cuts with the same state preserves endpoints
and visited states; the difficulty is always the channels, where a deletion
can bring two cuts into a comparison that did not exist before.

\begin{definition}[Fresh words]\label{def:fresh}
Coordinate $i$ of a tuple $W$ is \emph{fresh} in a table if its endpoint is
not visited earlier: $q_i(a)\ne q_i(\len{w_i})$ for every $a<\len{w_i}$. A
tuple is \emph{fresh} if all its coordinates are. A tuple $U$ \emph{freshly
dominates} $W$ if $U$ is fresh and $U\dominates W$. A fresh tuple is
\emph{minimal} if no tuple that freshly dominates it has a smaller total
length.
\end{definition}

The empty word is fresh. Since domination is transitive, every fresh tuple is
freshly dominated by a minimal fresh tuple: a fresh dominator of least total
length.

\begin{definition}[Immortal states]\label{def:immortal}
A state $q$ of a structural reader is \emph{immortal} if for every $m$ some
word of length $m$ can be read from $q$ without reaching $\dead$, and
\emph{mortal} otherwise.
\end{definition}

A state is immortal exactly when a cycle of states other than $\dead$ is
reachable from it without passing through $\dead$. Freshness comes for free
when the endpoint is mortal, and it yields the basic counting principle of
the whole proof: \emph{cuts at which comparable rests begin carry different
states}.

\begin{lemma}[Distinct states]\label{lem:distinct}
Let $W$ be a tuple and $a<b$ cuts of $w_i$ such that
$\wsuf{w_i}b\pre\wsuf{w_i}a$.
\begin{enumerate}
\item\label{it:distinct-fresh} If coordinate $i$ is fresh, then
  $q_i(a)\ne q_i(b)$.
\item\label{it:distinct-mortal} If coordinate $i$ is eligible with a mortal
  endpoint, then it is fresh.
\item\label{it:distinct-head} If coordinate $i$ is eligible and
  $b\le J_i$, then $q_i(a)\ne q_i(b)$.
\end{enumerate}
\end{lemma}

\begin{proof}
(1) If $q_i(a)=q_i(b)$, reading $\wsuf{w_i}b$ from the cut $a$ leads to the
endpoint, at the cut $a+\len{w_i}-b<\len{w_i}$. (2) If the endpoint occurs at
a cut $a<\len{w_i}$, the run from $a$ to the end is a cycle of states other
than $\dead$, so the endpoint is immortal. (3) If $q_i(a)=q_i(b)$, reading
$\wsub{w_i}b{J_i}$ from the cut $a$ leads to the rejected state $q_i(J_i)$ at
the cut $a+J_i-b<J_i$, before the first rejected cut.
\end{proof}

In particular, a fixed text has fewer than $N$ nonempty rests of one fresh
word as prefixes, since these rests are comparable.

\subsection{The Two Bounds}\label{sec:overview-bounds}

The first bound concerns the coordinates with immortal endpoints. Recall that
the rejected tail of an eligible coordinate is the part of its word after
the first rejected cut.

\begin{theorem}[Rejected tails]\label{thm:tails}
There is a computable function $\TailB$ such that every structural reader
with $k$ starts and $N$ states that has a hole has a hole in which the
rejected tail of every coordinate with an immortal endpoint has length at
most $\TailB(k,N)$.
\end{theorem}

The second bound is a statement about tables, which have no rejected states,
no dead state and no admitted pairs.

\begin{theorem}[Fresh tuples]\label{thm:fresh}
There is a computable function $\FreshB$ such that every minimal fresh tuple
of a table with $k$ starts and $N$ states has total length at most
$\FreshB(k,N)$. Hence every fresh tuple is freshly dominated by a tuple of
total length at most $\FreshB(k,N)$.
\end{theorem}

Neither bound depends on the size of the alphabet. \Cref{thm:fresh} is proved
in \cref{sec:fresh} and \cref{thm:tails} in \cref{sec:tails}.

Why two bounds? A hole need not be fresh: a coordinate with an immortal
endpoint may run through its endpoint many times. Deleting such a loop is
tempting, but it changes the rests of the word after all earlier cuts, and
with them the channels in which this coordinate is the host. For a
coordinate with a mortal endpoint the situation is rigid, by
\cref{lem:distinct}. \Cref{thm:tails} is the tool that brings immortal
coordinates under control: once the rejected tail is short, the coordinate
can be made fresh in a larger table that counts the letters of the tail.

\subsection{The Cover of a Hole}\label{sec:overview-cover}

%% cover -- the cover of a hole on a tracked coordinate.
\begin{figure}[t]
\centering
\begin{tikzpicture}[font=\scriptsize, >={Stealth[length=1.8mm]}, x=1.45cm, y=1cm,
    st/.style={draw, circle, minimum size=6.5mm, inner sep=0pt},
    rj/.style={st, double, double distance=1pt},
    cv/.style={draw, rounded corners=2pt, minimum height=6mm, minimum width=9mm, inner sep=1.5pt}]
%% the run in M
\node[st] (m0) at (0,0) {$s_i$};
\node[st] (m1) at (1,0) {$q$};
\node[st] (m2) at (2,0) {$q'$};
\node[rj] (m3) at (3.3,0) {$r_i$};
\node[rj] (m4) at (5.0,0) {$r'$};
\draw[->] (m0) -- node[above] {$c_1$} (m1);
\draw[->] (m1) -- node[above] {$c_2$} (m2);
\draw[->] (m2) -- node[above] {$c_3$} (m3);
\draw[->] (m3) to[bend left=22] node[above] {$c_4,\,c_6$} (m4);
\draw[->] (m4) to[bend left=22] node[below] {$c_5,\,c_7$} (m3);
\node[left] at (-0.45,0) {in $M$:};
%% the run in the cover
\node[cv] (c0) at (0,-1.7) {$(i,s_i)$};
\node[cv] (c1) at (1,-1.7) {$(i,q)$};
\node[cv] (c2) at (2,-1.7) {$(i,q')$};
\node[cv] (c3) at (3.3,-1.7) {$(i,0)$};
\node[cv] (c4) at (4.3,-1.7) {$(i,1)$};
\node[cv] (c5) at (5.3,-1.7) {$(i,2)$};
\node[cv] (c6) at (6.3,-1.7) {$(i,3)$};
\node[cv] (c7) at (7.3,-1.7) {$(i,4)$};
\draw[->] (c0) -- node[above] {$c_1$} (c1);
\draw[->] (c1) -- node[above] {$c_2$} (c2);
\draw[->] (c2) -- node[above] {$c_3$} (c3);
\draw[->] (c3) -- node[above] {$c_4$} (c4);
\draw[->] (c4) -- node[above] {$c_5$} (c5);
\draw[->] (c5) -- node[above] {$c_6$} (c6);
\draw[->] (c6) -- node[above] {$c_7$} (c7);
\node[left] at (-0.45,-1.7) {in $\cov M$:};
%% projection
\node[left] at (-0.45,-2.45) {$\proj$:};
\foreach \x/\l in {0/s_i,1/q,2/q',3.3/r_i,4.3/r',5.3/r_i,6.3/r',7.3/r_i} { \node at (\x,-2.45) {$\l$}; }
\foreach \x in {0,1,2,3.3,4.3,5.3,6.3,7.3} { \draw[densely dotted, ->] (\x,-2.02) -- (\x,-2.3); }
\draw[decorate, decoration={brace, mirror, amplitude=4pt}] (-0.3,-2.75) -- (2.3,-2.75)
  node[midway, below=4pt] {head: the states of $M$};
\draw[decorate, decoration={brace, mirror, amplitude=4pt}] (2.95,-2.75) -- (7.65,-2.75)
  node[midway, below=4pt] {rejected tail $c_4c_5c_6c_7$: one counter for each cut};
\end{tikzpicture}
\caption{The cover on a tracked coordinate $i$ with $w_i=c_1\cdots c_7$,
$J_i=3$ and $\ell_i=4$. In $M$ the rejected tail runs twice through the
cycle of the rejected states $r_i$ and $r'$, so the endpoint $r_i$ is
visited three times and the word is not fresh. In the cover the tail is
spelled out by the counters $(i,0),\dots,(i,4)$, the endpoint $(i,4)$ is
visited once. From each displayed counter, every transition other than the
displayed next-letter transition leads to the absorbing state $o$, and all
transitions from $(i,4)$ lead to $o$; other transitions from copies of head
states follow \cref{def:cover-hole}. On the displayed run, the projection $\proj$
(last row) maps each state to the state of $M$ at the same cut.}
\label{fig:cover}
\end{figure}


A hole whose immortal tails are short is not yet fresh, but it becomes fresh
in a larger table that counts the letters of these tails
(\cref{fig:cover}). Let $W$ be a hole of $M$ and $L$ a number such that
every coordinate with an immortal endpoint has a rejected tail of length at
most $L$. Call a coordinate $i$ \emph{tracked} if its rejected tail has
length $\ell_i=\len{w_i}-J_i\le L$; the other coordinates have mortal
endpoints.

\begin{definition}[Cover of a hole]\label{def:cover-hole}
The \emph{cover} $\cov M$ of $W$ is the following table with $k$ starts. Its
states are an absorbing state $o$ and, for each coordinate $i$, the pairs
$(i,s)$ with $s\in S$ and the \emph{counters} $(i,m)$ with $0\le m\le L$,
all of them distinct.
\begin{itemize}
\item For an untracked coordinate $i$, the start is $(i,s_i)$ and a letter
  $c$ takes $(i,s)$ to $(i,\delta(s,c))$.
\item For a tracked coordinate $i$ with first rejected state
  $r_i=q_i(J_i)$, a letter $c$ takes $(i,s)$ with $s\notin\Rej$ to the
  counter $(i,0)$ if $\delta(s,c)=r_i$, to $(i,\delta(s,c))$ if
  $\delta(s,c)\notin\Rej$, and to $o$ otherwise. It takes a counter $(i,m)$
  with $m<\ell_i$ to $(i,m+1)$ if $c$ is the letter of $w_i$ at the
  position $J_i+m$, and to $o$ otherwise; since $\ell_i\le L$, the counter
  $(i,m+1)$ exists. The start is $(i,0)$ if $J_i=0$ and $(i,s_i)$
  otherwise.
\item All remaining transitions lead to $o$. In particular, so do all
  transitions from the counter $(i,\ell_i)$ of a tracked coordinate and
  from the counters of an untracked one.
\end{itemize}
The \emph{projection} $\proj$ maps $(i,s)$ to $s$, the counter $(i,m)$ of a
tracked coordinate with $m\le\ell_i$ to $q_i(J_i+m)$, and all other counters
and $o$ to $\dead$.
\end{definition}

So the cover follows $M$ on the head of a tracked word and then spells out
its rejected tail, letter by letter. It has $k(N+L+1)+1$ states.

\begin{lemma}[Cover]\label{lem:cover}
The tuple $W$ is fresh in its cover $\cov M$, and every tuple that dominates
$W$ in $\cov M$ is a hole of $M$ with the same endpoints as $W$.
\end{lemma}

\begin{proof}
On a tracked coordinate $i$, the run of $w_i$ in the cover is $(i,q_i(a))$
at the cuts $a<J_i$ and the counter $(i,a-J_i)$ at the cuts $a\ge J_i$; on an
untracked coordinate it is $(i,q_i(a))$. In the first case the endpoint is
the counter $(i,\ell_i)$, which occurs only at the last cut. In the
second case the endpoint of $w_i$ in $M$ is mortal, so coordinate $i$ is
fresh in $M$ by \cref{lem:distinct}, hence in the cover. So $W$ is fresh in
the cover, and its runs avoid $o$.

Whenever a transition of the cover from a state $z$ on a letter $c$ does not
lead to $o$, it leads to a state with projection $\delta(\proj(z),c)$; this is
checked case by case, and for a counter it holds because the letter read is
the letter of $w_i$. The starts project to the starts of $M$. Now let $U$
dominate $W$ in the cover. The runs of $U$ in the cover visit only states
visited by $W$, so they avoid $o$, and therefore the run of $U$ in $M$ is the
projection of its run in the cover. Applying $\proj$ to the three conditions of
domination in the cover gives $U\dominates W$ in the table of $M$, because
the run of $W$ in $M$ is also the projection of its run in the cover. By
\cref{lem:hole-dominates}, $U$ is a hole of $M$ with the endpoints of $W$.
\end{proof}

\begin{proof}[Proof of \cref{thm:small-hole} from
\cref{thm:tails,thm:fresh}]
Let $M$ have a hole. By \cref{thm:tails} it has a hole $W$ whose coordinates
with immortal endpoints have rejected tails of length at most
$L=\TailB(k,N)$. By \cref{lem:cover}, $W$ is fresh in its cover, a table with
$k$ starts and $k(N+L+1)+1$ states, so by \cref{thm:fresh} it is freshly
dominated there by a tuple $U$ of total length at most
\begin{equation}\label{eq:hole-bound}
  \Hole(k,N)=\FreshB\bigl(k,\;k\,(N+\TailB(k,N)+1)+1\bigr).
\end{equation}
By \cref{lem:cover} again, $U$ is a hole of $M$.
\end{proof}

The composition explains the shape of the final bound: the fresh-tuple bound
is applied to a table whose size is the rejected-tail bound. For each fixed
number $k$ of starts, both are doubly exponential in $N$
(\cref{thm:elementary}), which gives four iterated exponentials in $N$.
